% !TEX program = pdflatex
\documentclass[12pt,a4paper]{letter}

\usepackage[T1]{fontenc}
\usepackage[margin=1in]{geometry}
\usepackage{hyperref}
\usepackage{setspace}
\usepackage{parskip}

\hypersetup{
  colorlinks=true,
  urlcolor=blue
}

\onehalfspacing

% Push closing/signature block to the right margin
\longindentation=0.65\textwidth

\signature{Professor {[FULL NAME]}\\
Professor of Computer Science\\
{[Department / Faculty]}\\
{[University Name]}, {[Country]}}

\address{{[Department of Computer Science]}\\
{[University Name]}\\
{[City, State/Province]}\\
{[Country]}}

\date{27 February 2026}

\begin{document}

\begin{letter}{UK Visas \& Immigration\\
Global Talent Visa --- Tech Nation (Digital Technology)\\
United Kingdom}

\opening{Dear Sir/Madam,}

I am writing to endorse Odeyemi Olusegun Israel for the UK Global Talent Visa in Digital Technology. I seldom write these letters. When I do, it is because someone's work has made me stop and look more carefully. Olusegun's did.

I am a Professor of Computer Science at {[University Name]}, where I work on {[relevant areas, e.g., distributed systems, machine learning, software engineering]}. I have known Olusegun as a mentor for roughly 10+ years---long enough to have watched him grow from a capable engineer into something harder to categorise: part systems builder, part theorist, and answerable mostly to his own curiosity.

\textbf{On the engineering work.} Olusegun has built a platform called AMTTP---the Anti-Money Laundering Transaction Trust Protocol. I will be blunt: it is more ambitious than most projects I see from funded research groups, and I include some of my own in that comparison. The system spans four architectural layers: a client-facing SDK layer (TypeScript and Python, both open-sourced), a compliance orchestration engine, an offline ML training pipeline, and a blockchain smart contract layer with 18 contracts deployed on Ethereum Sepolia. Seventeen containerised microservices sit behind a central gateway and have to talk to each other reliably.

That list, to someone outside the field, might look like a collection of buzzwords. It is not. Making those components work together---at scale, with deterministic compliance guarantees---is genuinely difficult. I have watched postdoctoral researchers struggle with far simpler integration tasks. Olusegun carried this largely by himself, which I find remarkable. It is one of the main reasons I agreed to write this letter.

The core technical insight is worth explaining briefly. In traditional finance, a suspicious transaction can be halted mid-flight. On a blockchain, once a transaction reaches finality, it is irreversible. Every existing AML tool---Chainalysis, Elliptic, and their competitors---operates after this point: they flag, but they do not prevent. AMTTP operates \textit{before} settlement, combining ML risk scores, graph topology analysis, sanctions list screening, and configurable policy rules into a single pass/fail decision. I have surveyed the literature, and as far as I can tell, no other system does this at the protocol level. The architectural contribution is real.

\textbf{On the machine learning and systems research contributions.} There are two lines of work I want to highlight. The one I find more striking now is aimed at critical infrastructure, not finance---though both matter.

The first is his paper \textit{Blind Spot Decomposition Theory: A Mathematical Framework for Characterising and Correcting Missed Fraud in Machine Learning Systems}\footnote{Available at \url{https://doi.org/10.5281/zenodo.18870589}.} (BSDT). It started as a way of explaining why ML systems miss the cases they miss, and has since grown into a broader framework for detecting structural instability in complex systems. The core idea: do not just flag that a model is failing, but decompose the failure into interpretable components and use those to guide correction. In the original formulation, Olusegun showed that missed detections break down into four parts---Camouflage, Feature Gap, Activity Anomaly, and Temporal Novelty---and this decomposition held across model families and domains without retraining or access to model internals. The rigour of that initial result is what first caught my attention.

But the recent work goes further. In his paper \textit{Blind-Spot Decomposition and the Geometry of System Collapse}\footnote{Available at \url{https://doi.org/10.5281/zenodo.19038783}.}, he proves that the BSDT energy functional is topologically equivalent to the interaction potential of a coupled dynamical system (Theorem~C: Morse-index equivalence at every critical point, homotopy equivalence of sublevel sets, two-sided gradient bounds). He then derives a spectral phase-transition theorem showing that below a critical manifold, any constant coupling stabilises the system, but above it, no fixed parameter works---only adaptive friction succeeds. He applies this framework to ERCOT power-grid dispatch data and achieves an AUC of 0.9997 with a 72-hour lead time before the rolling blackouts of Winter Storm Uri, with adaptive friction reducing peak capacity loss by 71\%. I take this seriously for a straightforward reason: a good result in financial prediction is commercially useful; a good result in power-grid monitoring has direct implications for safety and public infrastructure. Olusegun is no longer just building tools for economic diagnostics. He is building tools that can anticipate failure in systems people depend on daily, and that shift in scope is what I find most significant about his trajectory.

The same paper validates on FDIC banking data (6-quarter lead on the Global Financial Crisis, AUROC of 0.867, zero false alarms), the Terra/Luna crypto collapse ($r = 0.996$ correlation, 30-hour advance warning), and extends to three-dimensional incompressible Navier--Stokes equations with a 34-run GPU stress test at resolutions up to $256^3$ and Reynolds numbers up to 62,832. Nine contributions, three entirely separate physical domains, and a single underlying geometric structure connecting them.

The second line is his paper \textit{Universal Detection Principle: A Multi-Spectrum Tensor Framework for Domain-Agnostic Anomaly Detection}\footnote{Available at \url{https://doi.org/10.5281/zenodo.19037687}.} (UDP), which comes at anomaly detection from the other direction. Instead of explaining why a detector fails after the fact, UDP provides a tensor-based measure of deviation across multiple representations simultaneously. Anomalies that single-spectrum approaches miss can be surfaced in a principled way. The key point is not that it works on one benchmark---most methods do---but that it generalises.

These two threads are more connected than they might first appear. One explains failure structurally; the other builds a broader detection foundation. And his latest results show the ideas travel from finance into energy systems and fluid dynamics. I have supervised enough doctoral students to know that this kind of conceptual coherence across projects does not happen by accident. It comes from someone who is actually thinking, not just producing.

\textbf{On the broader system design.} A few details in the architecture are worth singling out, because they reveal a kind of engineering judgment you do not always see.

He has implemented a zero-knowledge proof framework (zkNAF) that lets institutions verify KYC credentials, risk score ranges, and sanctions non-membership without exposing personal data. The cryptographic engineering involved is non-trivial. More to the point, it sits right in the middle of an unresolved regulatory tension---GDPR privacy obligations on one side, AML transparency requirements on the other---and his design actually navigates it. He also built multi-oracle threshold signatures and replay protection, which is the kind of infrastructure-level security work that never makes it into a paper abstract but keeps real systems from falling apart.

The ML training pipeline is worth a word as well. It uses a Composite Teacher architecture: Autoencoder-Enhanced XGBoost (weighted at 0.4), seven FATF-aligned AML pattern detectors (0.3), and graph structural properties (0.3), generating pseudo-labels across 2.64 million transactions. A Student model then learns from those labels. Teacher-student distillation is the right call in a domain where ground-truth labels are scarce and expensive. The weighting choices themselves tell me he has thought carefully about how much to trust each signal source---something a less experienced engineer would simply leave at uniform.

\textbf{My assessment.} Over the years I have supervised and examined a fair number of PhD theses and postdoctoral projects. Some of them, if I am being candid, contained less original thinking than what Olusegun has produced on his own, outside of any university. I am not entirely comfortable saying that, but it is true.

His work sits across machine learning, distributed systems, blockchain architecture, and financial regulation. Each of those is a deep field. He has produced functioning software, original theory, and experimentally validated results in all four---not as part of a large team, and not with grant funding or institutional backing. I do not know many people who could do that. The ones I can think of all ended up running their own research programmes eventually.

Fintech alone would make him relevant to the UK, which leads Europe in that sector by most reasonable measures. But I would not characterise his contribution so narrowly. The work on power-grid stability, Navier--Stokes extension, and cross-domain transfer tells me he can build digital systems and mathematical tools that matter well beyond finance---in energy, in infrastructure resilience, in any setting where getting a prediction wrong has real consequences. The UK will need people who can connect secure engineering with machine learning and trustworthy decision-support in those domains. Olusegun is one of them. I recommend his application with confidence.

\vspace{1.5\baselineskip}

\closing{Yours faithfully,}

\end{letter}

\end{document}
